
\subsubsection{2.1 Reinforcement Learning from Execution Feedback for Code}

The use of execution signals to improve code language models is established in the literature. CodeRL \citep{Le2022CodeRL} used binary pass/fail rewards in an actor-critic framework, reporting improvements over SFT on APPS and HumanEval (+4.3\% pass@1). PPOCoder \citep{Majeed2023} applied PPO with binary execution reward, achieving 5--8\% gains on HumanEval and MBPP. RLTF \citep{Liu2023RLTF} introduced fractional test-passing reward and reported approximately 7\% improvement on HumanEval and MBPP, with the claim that partial credit reward advantages are especially pronounced at harder problem instances.

A common limitation of these works is that evaluation is confined to HumanEval and MBPP, covering only easy-to-medium difficulty. The question of whether RLEF advantages scale into hard competitive programming regimes remains unaddressed by these studies.

The most directly related prior study is RLEF \citep{Gehring2024RLEF}, which reports that partial reward outperforms binary reward more strongly at harder problems. Their study uses a proprietary Meta internal model, making independent replication impossible. The present work provides a reproducible counterpart using fully public infrastructure.

\subsubsection{2.2 Reward Formulation in RLEF}

Recent work has questioned whether fractional reward's denser signal produces better outcomes than binary reward. Two papers are directly relevant. First, a study at arXiv:2605.02944 finds that 97\% of APPS problems produce identical solve/fail outcomes regardless of reward formulation at GRPO convergence. Second, VeRPO (arXiv:2601.03525) identifies a cardinality bias in naïve fractional reward on variable-test-count datasets such as APPS: problems with single test cases make fraction reward identical to binary reward, and easy tests dominate the gradient signal without weighting corrections.

The null result reported in this paper (p=0.552) is consistent with both of these papers on a different model and dataset combination. The DeepSeek-R1 line of work \citep{DeepSeekAI2025R1} independently corroborates that binary execution feedback suffices at scale in mathematical reasoning domains.

\subsubsection{2.3 Hard Benchmark Evaluation for Code Language Models}

LiveCodeBench \citep{Jain2024LiveCodeBench} provides contamination-free evaluation stratified by difficulty, derived from competitive programming contests. MBPP \citep{Austin2021MBPP} and HumanEval \citep{Chen2021HumanEval} are established benchmarks covering easy and medium function-level problems respectively. APPS \citep{Hendrycks2021APPS} is the most widely used algorithmic training dataset and includes problems at introductory, interview, and competition difficulty. The coverage analysis in this paper --- 85.32\% of APPS competition problems have reference solutions --- is a novel empirical characterization that has not been reported in prior work.

\subsubsection{2.4 Positioning}

No prior work provides a controlled, reproducible, open-source comparison of RLEF versus SFT across the full difficulty spectrum from easy (HumanEval) to hard competitive programming (LiveCodeBench-Hard). The present work occupies the reproducibility gap between non-public large-scale studies (RLEF-2024, DeepSeek-R1) and easy-to-medium evaluations (CodeRL, PPOCoder, RLTF).

